% =============================================================
%  Kapitel 5 – Detailed Design: Emergency Stop
% =============================================================
\section{Detailed Design: Emergency Stop}
\label{sec:detaileddesign}

Dieser Abschnitt fokussiert ausschließlich auf das
\texttt{EmergencyStop}-Modul und beschreibt sein Klassendesign,
den internen Ablauf sowie den vollständigen Zustandsautomaten.

% ------------------------------------------------------------
\subsection{Klassendesign}
\label{sec:estop_klassen}

\begin{figure}[H]
  \centering
  \includegraphics[width=0.95\textwidth]{06_EStop_Class_Detailed}
  \caption{Detailliertes Klassendiagramm – EmergencyStop inkl.\ private Methoden}
  \label{fig:estop_class}
\end{figure}

Die Klasse \texttt{EmergencyStop} kapselt den vollständigen Sicherheitszustand.

\subsubsection{Attribute}

\begin{center}
\renewcommand{\arraystretch}{1.3}
\begin{tabular}{>{\ttfamily}p{4.5cm} >{\ttfamily}p{3cm} p{6cm}}
  \toprule
  \normalfont\textbf{Name} & \normalfont\textbf{Typ} & \normalfont\textbf{Beschreibung} \\
  \midrule
  currentState   & EStopState  & Aktueller Zustand des Automaten \\
  errorFlag      & bool        & Aktiver Fehler vorhanden \\
  safetyZoneActive & bool      & Safety-Zone (Rot) gerade aktiv \\
  driveRef       & DriveControlUnit* & Referenz auf Motorsteuerung \\
  loggerRef      & Logger*     & Referenz auf Logging-Modul \\
  \bottomrule
\end{tabular}
\end{center}

\subsubsection{Öffentliche Operationen (Public Interface)}

\begin{center}
\renewcommand{\arraystretch}{1.3}
\begin{tabular}{>{\ttfamily}p{6.5cm} p{7cm}}
  \toprule
  \normalfont\textbf{Signatur} & \normalfont\textbf{Beschreibung} \\
  \midrule
  onZoneUpdate(zone: Zone) : void
    & Eingehende Zonenänderung verarbeiten (Safety / Normal). \\
  onEmergencyButton() : void
    & Pilzknopf-Event verarbeiten. \\
  onFrontTrigger() : void
    & Fronttaster-Event verarbeiten. \\
  onReleaseRequest() : void
    & Manuelle Freigabe anfordern; Guard wird geprüft. \\
  onErrorDetected() : void
    & Fehler-Ereignis (z.\,B.\ fehlende Sensordaten) verarbeiten. \\
  onShutdownRequest() : void
    & Kontrollierter Shutdown einleiten. \\
  \bottomrule
\end{tabular}
\end{center}

\subsubsection{Private Methoden}

Alle internen Handler sind als \textbf{private Methoden} implementiert,
um die Kapselung der Zustandslogik zu gewährleisten.

\begin{center}
\renewcommand{\arraystretch}{1.3}
\begin{tabular}{>{\ttfamily}p{6.5cm} p{7cm}}
  \toprule
  \normalfont\textbf{Signatur} & \normalfont\textbf{Beschreibung} \\
  \midrule
  triggerEmergencyStop() : void
    & Setzt Zustand auf \texttt{EmergencyStop}; ruft \texttt{stopDrive()} auf. \\
  triggerSoftStop() : void
    & Setzt Zustand auf \texttt{SoftStop}; ruft \texttt{stopDrive()} auf. \\
  stopDrive() : void
    & Sendet Stop-Befehl an \texttt{DriveControlUnit}. \\
  enableDrive() : void
    & Sendet Freigabe-Befehl an \texttt{DriveControlUnit}. \\
  checkReleaseGuard() : bool
    & Prüft: \texttt{!safetyZoneActive \&\& !errorFlag}. \\
  logEvent(level: LogLevel, msg: string) : void
    & Delegiert an \texttt{Logger}; fügt Zeitstempel hinzu. \\
  transitionTo(state: EStopState) : void
    & Führt Zustandsübergang durch und protokolliert ihn. \\
  \bottomrule
\end{tabular}
\end{center}

\subsubsection{Enum: EStopState}

\begin{center}
\renewcommand{\arraystretch}{1.2}
\begin{tabular}{>{\ttfamily}p{4cm} p{9cm}}
  \toprule
  \normalfont\textbf{Wert} & \normalfont\textbf{Bedeutung} \\
  \midrule
  Initialization & Startphase; alle Aktuatoren gesperrt. \\
  SoftStop       & Sicherer Stopp; automatisch wiederherstellbar nach Freigabe. \\
  Normal         & Normalbetrieb; \texttt{DriveControlUnit} freigegeben. \\
  EmergencyStop  & Notfall-Stopp; blockiert bis manuelle Freigabe + Guard-Check. \\
  HardStop       & Höchste Priorität; hardware-seitiger Stopp (kein SW-Override). \\
  \bottomrule
\end{tabular}
\end{center}

\todobox{Vollständige Liste aller privaten Methoden mit Vor-/Nachbedingungen
(Design by Contract) noch auszuarbeiten.}

% ------------------------------------------------------------
\subsection{Interner Ablauf (Sequenzdiagramm)}
\label{sec:estop_seq}

\begin{figure}[H]
  \centering
  \includegraphics[width=0.95\textwidth]{06_EStop_Internal_Sequence}
  \caption{Internes Sequenzdiagramm – Ablauf innerhalb von EmergencyStop}
  \label{fig:estop_seq}
\end{figure}

Das interne Sequenzdiagramm (Abbildung~\ref{fig:estop_seq}) zeigt,
wie ein eingehender Event durch die privaten Methoden verarbeitet wird.
Die Verarbeitung folgt immer dem gleichen Schema:

\begin{enumerate}
  \item \textbf{Event-Empfang} über eine öffentliche Methode
        (z.\,B.\ \texttt{onZoneUpdate()}).
  \item \textbf{Guard-Prüfung} (falls erforderlich) durch \texttt{checkReleaseGuard()}.
  \item \textbf{Zustandsübergang} über \texttt{transitionTo()} mit automatischem Logging.
  \item \textbf{Aktorsteuerung}: Aufruf von \texttt{stopDrive()} oder \texttt{enableDrive()}.
  \item \textbf{Logging} über \texttt{logEvent()} mit Level und Zeitstempel.
\end{enumerate}

Diese Trennung hält die Event-Handler frei von direkten Hardware-Zugriffen
und ermöglicht einfaches Unit-Testing.

% ------------------------------------------------------------
\subsection{Zustandsautomat (State Machine)}
\label{sec:estop_statemachine}

\begin{figure}[H]
  \centering
  \includegraphics[width=0.95\textwidth]{06_EStop_StateMachine}
  \caption{Zustandsautomat – vollständiges Verhalten von EmergencyStop}
  \label{fig:estop_statemachine}
\end{figure}

Der Zustandsautomat (Abbildung~\ref{fig:estop_statemachine}) definiert das
vollständige Verhalten von \texttt{EmergencyStop}.

\subsubsection{Zustandsübergänge}

\begin{center}
\renewcommand{\arraystretch}{1.3}
\begin{tabular}{>{\ttfamily}p{3.5cm} p{4cm} >{\ttfamily}p{3.5cm}}
  \toprule
  \normalfont\textbf{Von} & \normalfont\textbf{Trigger / Guard} & \normalfont\textbf{Nach} \\
  \midrule
  Initialization    & SystemReady                          & SoftStop \\
  SoftStop          & ReleaseRequest [Guard OK]            & Normal \\
  Normal            & ZoneUpdate(Safety)                   & EmergencyStop \\
  Normal            & EmergencyButton / FrontTrigger       & EmergencyStop \\
  Normal            & ErrorDetected                        & SoftStop \\
  Normal            & ShutdownRequest                      & SoftStop \\
  EmergencyStop     & ReleaseRequest [Guard OK]            & Normal \\
  EmergencyStop     & ReleaseRequest [Guard FAIL]          & EmergencyStop \\
  EmergencyStop     & ErrorDetected                        & SoftStop \\
  SoftStop          & ErrorDetected                        & SoftStop \\
  * (alle)          & HardStopSignal                       & HardStop \\
  \bottomrule
\end{tabular}
\end{center}

\subsubsection{Besondere Eigenschaften}

\begin{itemize}
  \item \textbf{HardStop} hat die höchste Priorität; er kann aus jedem Zustand
        erreicht werden und ist nicht softwareseitig aufhebbar.
  \item \textbf{EmergencyStop} erfordert zwingend eine manuelle Freigabe plus
        erfolgreiche Guard-Prüfung.
  \item \textbf{ErrorDetected} führt immer zu \texttt{SoftStop}
        (Fail-Safe-Prinzip).
  \item \textbf{ShutdownRequest} führt immer zu \texttt{SoftStop}.
\end{itemize}

\todobox{Activity-Diagramme für die wichtigsten privaten Methoden
(\texttt{checkReleaseGuard}, \texttt{triggerEmergencyStop}) ausarbeiten
oder begründen, warum die State-Machine hinreichend ist.}
